\scanpage{268}
\margin{\figure{\begin{tikzpicture}\draw[<-] (0,0)--(.5,0);\end{tikzpicture}}}
- \mspage{83} -
\margin{IX}
7.1.3. La projectivité de \(\hat B\) signifie, en vertu de EGA III 5.4.5, que l'on peut trouver un Module inversible \(\hat{\uline{M}}\) sur \(\hat B,\) i.e. une suite de Modules inversibles \(\uline{M}_n\) sur les \(B_n\) qui se ``\unsure{recollent}'' pour \(n\) variable, qui soit \uline{ample} relativement à \(S\) i.e. qui soit tel que \(M_0\) soit ample. L'idée naturelle serait de partir d'un Module inversible ample \(\uline{M}\) de \(G,\) et de montrer que pour tout \(n,\) la restriction \(\uline{L}_n\) de \(\uline{L}\) à \(G_n\) provient par image inverse d'un Module inversible \(M_n\) sur \(B_n.\) J'ignore si de tels \(M_n\) existent nécessairement (et en doute fort). Nous allons suivre une méthode un peu différente, en construisant une \uline{polarisation} de \(B,\) i.e. un système projectif de \uline{polarisations} \(\rho_n\) des \(B_n\) qui se recollent. On sait alors [5] que les doubles de ces polarisations définissent canoniquement des Modules inversibles \(\uline{M}_n\) (symétriques) sur les \(B_n,\) qui par leur nature canonique se recolleront nécessairement, et fourniront le \(\unsure{\hat{\uline{M}}|\uline{M}}\) cherché.

Nous allons interpréter un élément de \(\operatorname{NS}(B_n/S_n)\) comme une correspondance divisorielle symétrique sur \((B_n,B_n)\) [5], i.e. comme une biextension de \((B_n,B_n)\) par \(\mathbb G_{m\,S_n}.\) En vertu de VIII 3.5, ces biextensions correspondent biunivoquement aux biextensions de \((A_n,A_n)\) par \(\mathbb G_{mS}.\) Pour obtenir un système cohérent de telles biextensions \(W_n\) il suffit de partir d'une
\[
(7.1.4)\qquad W,\ \text{biextension symétrique de }(A,A)\text{ par }\mathbb G_{mS},
\]
en prenant les \(W_n\) se déduisant de \(W\) par les changements de base \(S_n\longrightarrow S.\) Nous devons de plus supposer que l'élément de \(\operatorname{NS}(B_0)\) défini par \(W_0\) est une polarisation de \(B_0\) pour obtenir ainsi une polarisation de \(\hat B.\)
\figure{\begin{tikzpicture}\fill (0,0)--(.06,.10)--(.12,.02)--(.20,.13)--(.16,-.03)--(.06,-.04)--cycle;\end{tikzpicture}}
\margin{401}
\margin{267}
\margin{401}
